\section{Uniform profile overlap on sampled orbit fibers}
\label{sec:tpc52-profile-overlap}

The diagonal Mellin identity reduces the actual-direction problem to a
nonnegative comparison: the sampled energy of the unseparated physical
profile must not be much smaller than the energy of the unimodular Mellin
carrier.  This section proves that comparison under one explicit geometric
hypothesis.  It is important that the hypothesis is stated separately from
the row-flatness theorem inherited from TPC--50.

\subsection{A row--orbit sampling interface}

Put \(T=\log 2\) and \(Y=[0,T]^2\).  Let \(\cR_X\) be a finite
opposite-row set, with logarithmic coordinates
\[
 q_\gamma=(x_\gamma,y_\gamma)\in Y
 \qquad(\gamma\in\cR_X),
\]
positive weights \(\omega_{\gamma,X}\), and total weight
\begin{equation}
 E_X:=\sum_{\gamma\in\cR_X}\omega_{\gamma,X}>0.
 \label{eq:tpc52-opposite-row-energy}
\end{equation}
We use the following three hypotheses.

\begin{enumerate}[label=\textup{(S\arabic*)},leftmargin=3.2em]
 \item The probability measures
 \begin{equation}
  \nu_X:=E_X^{-1}\sum_{\gamma\in\cR_X}
           \omega_{\gamma,X}\delta_{q_\gamma}
  \label{eq:tpc52-row-empirical-measure}
 \end{equation}
  converge weakly to a probability measure \(\nu\) on \(Y\).  In the
  prime--squarefree application one takes the exact base weight
  \(\omega_{\gamma,X}=(\log\ell_\gamma)^{2e}\), \(e\in\{0,1\}\), and
  places every fixed dyadic or unary shape in \(B\) and \(F\).  Then
  \(d\nu=e^{x+y}\,dx\,dy\) in the limit (this density already has mass
  one on \(Y\)) \citep{WangTPC51}.  With any other base weight, both
  weak convergence and the identity of its limiting measure are separate
  hypotheses.
 \item The base weights are uniformly flat:
 \begin{equation}
  |\cR_X|\max_{\gamma\in\cR_X}\omega_{\gamma,X}
  \le C_{\rm flat}E_X.
  \label{eq:tpc52-row-flatness}
 \end{equation}
 \item The literal mask \(M_X(\alpha,\gamma)\in\{0,1\}\) obeys the
 one-sided degree estimate
 \begin{equation}
  \sup_\alpha\#\{\gamma\in\cR_X:M_X(\alpha,\gamma)=0\}
  \le\epsilon_X|\cR_X|,
  \qquad \epsilon_X=o(1).
  \label{eq:tpc52-mask-degree}
 \end{equation}
\end{enumerate}

For the TPC--50/51 application, \(\cR_X\) is required to be one complete
positive-density admissible prime--squarefree row cell after fixed residue
subdivision, in the nondegenerate interior regime used there (in particular
the inherited size separation and injectivity of \(m=\ell d\) hold).  No
independent pruning of the opposite row is allowed before \textup{(S3)} is
applied.  If a declared packet does not have this form, \textup{(S1)}--
\textup{(S3)} are new assumptions rather than consequences of TPC--50/51.

The source label \(\alpha\) is not assigned a probability measure.  It may
be supported on one row, and later it may carry an arbitrary cofactor
multiplicity.  This asymmetry is what permits the final theorem to be
uniform under source and \(r\)-concentration.

Let \(\mathfrak B\) be a compact parameter space containing all normalized
source-scale parameters.  Let \(\cA_0\) be a fixed finite collection of
direct orbit-residue cells.  To each \(a\in\cA_0\) and
\(\beta\in\mathfrak B\) associate a carrier and a physical profile
\begin{equation}
 B_{\beta,a},F_{\beta,a}\in C^1(Y\times I_0).
 \label{eq:tpc52-carrier-profile-family}
\end{equation}
We assume that both families are compact in \(C^1(Y\times I_0)\).  In
particular, their sup norms and first derivatives are uniformly bounded.
The direct cell \(a\) corresponds to the progression
\(j\equiv a\pmod {H_0}\); cells with zero physical residue amplitude are
removed before this definition.

Define the continuous carrier and profile energies by
\begin{align}
 \mathscr B(\beta)
 &:=\frac1{H_0}\sum_{a\in\cA_0}
   \int_Y\int_{I_0}|B_{\beta,a}(q,t)|^2\,dt\,d\nu(q),
 \label{eq:tpc52-continuous-carrier}\\
 \mathscr P(\beta)
 &:=\frac1{H_0}\sum_{a\in\cA_0}
   \int_Y\int_{I_0}|F_{\beta,a}(q,t)|^2\,dt\,d\nu(q).
 \label{eq:tpc52-continuous-profile}
\end{align}
We impose the carrier nondegeneracy and normalization
\begin{equation}
 0<b_-\le \mathscr B(\beta)\le b_+<\infty
 \qquad(\beta\in\mathfrak B).
 \label{eq:tpc52-carrier-normalization}
\end{equation}
Deleting an identically zero direct cell is harmless, but the uniform lower
bound \(b_->0\) is a genuine compact-packet hypothesis; it is not supplied
by row flatness alone.  Once that lower bound is known, rescaling fixes the
upper and lower constants.

\begin{assumption}[Uniform physical profile nondegeneracy]
\label{ass:tpc52-profile-nondegeneracy}
The declared direct packet satisfies
\begin{equation}
 \boxed{
 \rho_*:=\inf_{\beta\in\mathfrak B}
          \frac{\mathscr P(\beta)}{\mathscr B(\beta)}>0.}
 \label{eq:tpc52-uniform-overlap}
\end{equation}
\end{assumption}

Assumption~\ref{ass:tpc52-profile-nondegeneracy} is new.  It is not a
consequence of TPC--50.  That paper proves flatness and literal-mask survival
for an identified fixed-complexity unary row factor, uniformly in one
real-axis Mellin phase.  It does not prove that the complete product of
source, opposite-row, and orbit cutoffs retains a positive amount of carrier
energy after the Mellin variables are reassembled.  In particular, it does
not exclude two scaled compact supports approaching a common flat endpoint.
For named fixed cutoffs, \eqref{eq:tpc52-uniform-overlap} is a deterministic
support-geometry condition on a compact finite-dimensional parameter set; it
is not a prime-pair hypothesis.

\begin{lemma}[A nonempty active subcell from one strict overlap]
\label{lem:tpc52-local-active-subcell}
Assume \(\mathfrak B\subset\R^d\), the maps
\(\beta\mapsto B_{\beta,a},F_{\beta,a}\) are continuous in
\(C(Y\times I_0)\), and \eqref{eq:tpc52-carrier-normalization} holds.
If, for some interior parameter \(\beta_0\), some direct cell \(a_0\),
and some \((q_0,t_0)\in\supp\nu\times\operatorname{int}I_0\),
\begin{equation}
 |F_{\beta_0,a_0}(q_0,t_0)|>0,
 \label{eq:tpc52-one-strict-overlap}
\end{equation}
then there is a nonempty compact parameter cell
\(\mathfrak B_0\Subset\mathfrak B\), containing \(\beta_0\) in its
relative interior, on which
\begin{equation}
 \inf_{\beta\in\mathfrak B_0}
 \frac{\mathscr P(\beta)}{\mathscr B(\beta)}>0.
 \label{eq:tpc52-local-active-ratio}
\end{equation}
\end{lemma}

\begin{proof}
Continuity and \eqref{eq:tpc52-one-strict-overlap} give an open
row--orbit box of positive \(\nu\otimes dt\) measure on which the same
profile remains nonzero.  Hence \(\mathscr P(\beta_0)>0\).  The two
energy functions are continuous in \(\beta\), while
\(\mathscr B\ge b_->0\).  Their ratio is therefore continuous and
positive on a sufficiently small neighborhood of \(\beta_0\).  Choose
a compact cell inside that neighborhood.
\end{proof}

The lemma proves existence only after a strict overlap point has been
verified for the named physical cutoffs.  It does not claim that every
dyadic boundary packet, or the closure of the full inherited parameter
range, is uniformly active.

For \(\beta\in\mathfrak B\), define the unmasked sampled energies
\begin{align}
 \mathscr B_X^0(\beta)
 &:=\frac1{E_XJ}\sum_{a\in\cA_0}
   \sum_{\gamma\in\cR_X}\omega_{\gamma,X}
   \sum_{j\in\cJ_{X,a}}
   |B_{\beta,a}(q_\gamma,j/J)|^2,
 \label{eq:tpc52-unmasked-carrier-sample}\\
 \mathscr P_X^0(\beta)
 &:=\frac1{E_XJ}\sum_{a\in\cA_0}
   \sum_{\gamma\in\cR_X}\omega_{\gamma,X}
   \sum_{j\in\cJ_{X,a}}
   |F_{\beta,a}(q_\gamma,j/J)|^2.
 \label{eq:tpc52-unmasked-profile-sample}
\end{align}
For a source row \(\alpha\), carrying a parameter
\(\beta_\alpha\in\mathfrak B\), insert
\(M_X(\alpha,\gamma)\) in these sums and denote the resulting quantities by
\(\mathscr B_X^M(\alpha)\) and \(\mathscr P_X^M(\alpha)\).

\begin{lemma}[Uniform row--orbit sampling]
\label{lem:tpc52-uniform-row-orbit-sampling}
Under \textup{(S1)} and the compact \(C^1\) profile hypotheses,
\begin{align}
 \sup_{\beta\in\mathfrak B}
 |\mathscr B_X^0(\beta)-\mathscr B(\beta)|&=o(1),
 \label{eq:tpc52-uniform-carrier-sampling}\\
 \sup_{\beta\in\mathfrak B}
 |\mathscr P_X^0(\beta)-\mathscr P(\beta)|&=o(1).
 \label{eq:tpc52-uniform-profile-sampling}
\end{align}
\end{lemma}

\begin{proof}
For a uniformly bounded subset of \(C^1(I_0)\), progression Riemann sums
obey
\[
 \frac1J\sum_{j\in\cJ_{X,a}}
 h(j/J)
 =\frac1{H_0}\int_{I_0}h(t)\,dt+O(J^{-1})
\]
uniformly in \(h\).  Apply this with
\(h(t)=|B_{\beta,a}(q,t)|^2\) and with \(F\) in place of \(B\).  The
resulting functions of \(q\), as \(\beta\) and \(a\) vary, form compact
subsets of \(C(Y)\).

Weak convergence of probability measures is uniform on a compact subset of
\(C(Y)\): given \(\delta>0\), choose a finite \(\delta\)-net in the sup norm,
use weak convergence on its finitely many elements, and then use that both
measures have mass one.  Applying this observation to the preceding compact
families, and then summing over the fixed set \(\cA_0\), proves both claims.
\end{proof}

\begin{lemma}[Literal-mask deletion]
\label{lem:tpc52-literal-mask-deletion}
Under \textup{(S2)}--\textup{(S3)}, uniformly in the source row \(\alpha\),
\begin{align}
 0\le
 \mathscr B_X^0(\beta_\alpha)-\mathscr B_X^M(\alpha)
 &\le C\epsilon_X,
 \label{eq:tpc52-carrier-mask-deletion}\\
 0\le
 \mathscr P_X^0(\beta_\alpha)-\mathscr P_X^M(\alpha)
 &\le C\epsilon_X,
 \label{eq:tpc52-profile-mask-deletion}
\end{align}
where \(C\) depends only on the fixed profile families, residue cells, and
the flatness constant.
\end{lemma}

\begin{proof}
For each deleted opposite row, the orbit sum of either squared profile is
\(O(J)\), uniformly in \(\alpha,\gamma\), by compactness in \(C^1\).
There are at most \(\epsilon_X|\cR_X|\) deleted rows.  Hence the unnormalized
loss is at most
\[
 C J\epsilon_X|\cR_X|
       \max_\gamma\omega_{\gamma,X}
 \le C C_{\rm flat}\epsilon_XJE_X
\]
by \eqref{eq:tpc52-row-flatness}.  Division by \(JE_X\) proves the two
estimates.
\end{proof}

\begin{theorem}[Uniform sampled overlap under the literal mask]
\label{thm:tpc52-uniform-sampled-overlap}
Assume \textup{(S1)}--\textup{(S3)},
\eqref{eq:tpc52-carrier-normalization}, and
Assumption~\ref{ass:tpc52-profile-nondegeneracy}.  Then, for all sufficiently
large \(X\), uniformly in every source row \(\alpha\),
\begin{equation}
 \mathscr B_X^M(\alpha)\ge \frac{b_-}{2}>0,
 \qquad
 \frac{\mathscr P_X^M(\alpha)}{\mathscr B_X^M(\alpha)}
 \ge \frac{\rho_*}{2}.
 \label{eq:tpc52-uniform-masked-overlap}
\end{equation}
\end{theorem}

\begin{proof}
By \cref{lem:tpc52-uniform-row-orbit-sampling,lem:tpc52-literal-mask-deletion},
both masked sampled energies differ
uniformly from their continuous counterparts by \(o(1)\).  In particular,
the carrier is at least \(b_-/2\) for large \(X\).  Also
\(\mathscr P(\beta)\ge\rho_*\mathscr B(\beta)\ge\rho_*b_-\).  Choose \(X\)
so large that the two absolute errors are smaller than a fixed sufficiently
small multiple of \(\rho_*b_-\).  Dividing the resulting profile lower bound
by the corresponding carrier upper bound gives the stated constant
\(\rho_*/2\), after harmlessly enlarging the threshold for \(X\).
\end{proof}

The number \(1/2\) is inessential.  The proof gives
\(\rho_*+o(1)\) after the carrier normalization is kept in the ratio.  The
weaker displayed constant is convenient for later direct sums.

\subsection{Arbitrary outer concentration}

Let \(\cU_X\) be any finite source/cofactor label set.  Each
\(u\in\cU_X\) carries a source row \(\alpha(u)\), a parameter
\(\beta_{\alpha(u)}\), and an arbitrary scalar or Hilbert coefficient
\(a_u\).  Define
\begin{align}
 \cC_X(a)&:=E_XJ\sum_{u\in\cU_X}
             \|a_u\|^2\mathscr B_X^M(\alpha(u)),
 \label{eq:tpc52-outer-carrier-energy}\\
 \cP_X(a)&:=E_XJ\sum_{u\in\cU_X}
             \|a_u\|^2\mathscr P_X^M(\alpha(u)).
 \label{eq:tpc52-outer-profile-energy}
\end{align}

\begin{corollary}[Source- and cofactor-uniform overlap]
\label{cor:tpc52-arbitrary-source-r-concentration}
Under the hypotheses of \cref{thm:tpc52-uniform-sampled-overlap},
\begin{equation}
 \boxed{\cP_X(a)\ge \frac{\rho_*}{2}\cC_X(a)}
 \label{eq:tpc52-outer-overlap}
\end{equation}
for every finitely supported family \((a_u)\).  No lower source count,
source equidistribution, or uniform cofactor multiplicity is required.
\end{corollary}

\begin{proof}
Apply \eqref{eq:tpc52-uniform-masked-overlap} separately to every retained
source/cofactor fiber, multiply by \(\|a_u\|^2\), and sum.  All terms are
nonnegative because the source and cofactor labels have not been coherently
grouped.
\end{proof}

\subsection{Good and bad overlap cells}

Uniform positivity may fail on the closure of a packet even though most
carrier energy lies away from the grazing boundary.  The following threshold
form records exactly what an energy selection must prove.

For \(\eta>0\), put
\begin{equation}
 \mathfrak B_{\rm good}(\eta)
 :=\left\{\beta\in\mathfrak B:
          \frac{\mathscr P(\beta)}{\mathscr B(\beta)}\ge\eta\right\},
 \qquad
 \mathfrak B_{\rm bad}(\eta):=\mathfrak B\setminus
          \mathfrak B_{\rm good}(\eta).
 \label{eq:tpc52-good-bad-parameters}
\end{equation}
Let \(\Delta_X^{\rm samp}\) denote the maximum of the two uniform sampling errors in
\cref{lem:tpc52-uniform-row-orbit-sampling} plus \(C\epsilon_X\), with the
constant from \cref{lem:tpc52-literal-mask-deletion}.

\begin{proposition}[Threshold overlap tradeoff]
\label{prop:tpc52-threshold-overlap}
Let \(\eta_X>0\) satisfy
\begin{equation}
 \Delta_X^{\rm samp}=o(\eta_X).
 \label{eq:tpc52-relative-sampling-threshold}
\end{equation}
For all sufficiently large \(X\), every source fiber whose parameter lies in
\(\mathfrak B_{\rm good}(\eta_X)\) has sampled masked profile-to-carrier
ratio at least \(\eta_X/2\).

For an arbitrary outer family \(a\), let \(\cC_X^{\rm good}(a)\) and
\(\cC_X^{\rm bad}(a)\) be the carrier energies from
\eqref{eq:tpc52-outer-carrier-energy} restricted according to
\eqref{eq:tpc52-good-bad-parameters}, and put
\begin{equation}
 \sigma_X:=\frac{\cC_X^{\rm good}(a)}{\cC_X(a)}
 \quad\text{when }\cC_X(a)>0.
 \label{eq:tpc52-good-carrier-fraction}
\end{equation}
Then
\begin{equation}
 \boxed{
 \cP_X(a)\ge \frac{\eta_X\sigma_X}{2}\cC_X(a).}
 \label{eq:tpc52-threshold-global-overlap}
\end{equation}
\end{proposition}

\begin{proof}
The proof of \cref{thm:tpc52-uniform-sampled-overlap} is uniform on the good
set.  Condition \eqref{eq:tpc52-relative-sampling-threshold} makes its
absolute sampling-and-deletion error negligible relative to
\(\eta_X\mathscr B(\beta)\), since the carrier is at least \(b_-\).
This gives the fiberwise factor \(\eta_X/2\).  Sum only the good fibers as in
\cref{cor:tpc52-arbitrary-source-r-concentration}; the bad contribution to
the physical energy is nonnegative and may be discarded.  The definition of
\(\sigma_X\) then gives \eqref{eq:tpc52-threshold-global-overlap}.
\end{proof}

Mere weak convergence in \textup{(S1)} proves the proposition for a fixed
\(\eta>0\), but it does not by itself verify
\eqref{eq:tpc52-relative-sampling-threshold} for a polynomially shrinking
\(\eta_X\).  Such a use requires a quantitative row-sampling theorem at the
resolution of the shrinking overlap, together with a mask error smaller than
\(\eta_X\).  This distinction prevents a qualitative compactness argument
from being misreported as a polynomial endpoint estimate.

For clarity, the quantitative fixed-modulus sampling error available in the
inherited row theorem contains \(e^{-c\sqrt{\log L}}\).  This tends to zero
but is larger than \(X^{-\lambda}\) for every fixed \(\lambda>0\).  Hence it
does not verify \(\Delta_X^{\rm samp}=o(X^{-\lambda})\) for a polynomially shrinking
threshold.  A strict inequality \(\lambda<1/400\) controls the inherited
mask-deletion term only; the shrinking-threshold branch still requires a
sharper relative sampling theorem (or a fixed positive overlap).

\subsection{The exponent ledger}

Suppose a displayed Mellin representation satisfies
\begin{equation}
 \cE_X^2\le X^{\omega+o(1)}\cC_X(a),
 \label{eq:tpc52-envelope-carrier-exponent}
\end{equation}
and write
\begin{equation}
 \eta_X=X^{-\lambda+o(1)},
 \qquad
 \sigma_X=X^{-\varsigma+o(1)}.
 \label{eq:tpc52-threshold-exponents}
\end{equation}
Then \cref{prop:tpc52-threshold-overlap} gives
\begin{equation}
 \frac{\cE_X^2}{\cP_X(a)}
 \le X^{\omega+\lambda+\varsigma+o(1)}.
 \label{eq:tpc52-threshold-condition-exponent}
\end{equation}
For a family of \(R_X=X^{\varrho+o(1)}\) direct cells whose individual
weighted Mellin masses are at most \(X^{\tau+o(1)}\), Cauchy--Schwarz in the
finite cell index gives the admissible ledger
\begin{equation}
 \omega\le 2\tau+\varrho.
 \label{eq:tpc52-cell-mass-exponent}
\end{equation}
Thus a sufficient endpoint condition is
\begin{equation}
 \boxed{2\tau+\varrho+\lambda+\varsigma\le \frac1{400},}
 \label{eq:tpc52-one-over-four-hundred-ledger}
\end{equation}
in addition to the relative sampling requirement
\eqref{eq:tpc52-relative-sampling-threshold}.  For a declared fixed collection
of direct residue cells, fixed smooth shapes, and fixed derivative order,
\(\tau=\varrho=0\).  Under the additional exact-profile identification and
uniform nondegeneracy hypotheses used below,
\(\lambda=\varsigma=0\) as well, so this gate consumes no positive part of
the \(1/400\) allowance.

At the inherited endpoint the degree defect contains a term
\(X^{-1/400+o(1)}\).  Consequently a threshold exactly of size
\(X^{-1/400+o(1)}\) is not automatically larger than the mask error.  A
strict exponent margin, or a sharper packet-specific deletion theorem, is
needed before equality at that boundary may be claimed.
